% 彗星（综述）
% license CCBYSA3
% type Wiki

本文根据 CC-BY-SA 协议转载翻译自维基百科\href{https://en.wikipedia.org/wiki/Comet}{相关文章}。

\begin{figure}[ht]
\centering
\includegraphics[width=8cm]{./figures/45ca94c34375617b.png}
\caption{} \label{fig_Comet_1}
\end{figure}
彗星是一类由冰组成的、体积较小的太阳系天体或星际天体。当其经过接近太阳的区域时，会因受热而开始释放气体，这一过程称为释气作用（outgassing）。这一过程会在彗核周围产生一个延展的、未被引力束缚的大气层，即彗发（coma），并有时形成由彗发中吹离的气体和尘埃构成的彗尾。这些现象源于太阳辐射以及向外流动的太阳风等离子体对彗核的作用。彗核的大小从数百米到数十公里不等，由松散聚集的冰、尘埃和小型岩石颗粒组成。彗发的尺度可达地球直径的 15 倍，而彗尾的长度甚至可能超过一个天文单位。如果距离足够近且亮度足够高，人类在无望远镜的情况下即可从地球上观测到彗星，其在天空中的角距可达 30°（相当于 60 个月亮）。自古以来，彗星已被众多文化和宗教观察并记录。

彗星通常具有高离心率的椭圆轨道，其轨道周期范围极广，从数年到可能长达数百万年。短周期彗星来源于海王星轨道之外的柯伊伯带或相关的散射盘。长周期彗星则被认为起源于奥尔特云，一个由冰质天体组成的球形云层，其范围从柯伊伯带外侧延伸至通往最近恒星方向一半的距离\(^\text{[2]}\)。长周期彗星通常因经过恒星的引力扰动或银河潮汐效应而被触发向太阳方向运动。双曲轨道彗星可能只会一次穿越内太阳系，随后被抛入星际空间。彗星的出现称为一次显现（apparition）。

多次接近太阳的灭绝彗星会失去几乎全部挥发性冰与尘埃，并可能逐渐呈现出类似小型小行星的外观\(^\text{[3]}\)。一般认为，小行星与彗星的起源不同：小行星形成于木星轨道以内，而非太阳系外部\(^\text{[4][5]}\)。然而，主带彗星（main-belt comets）以及活跃的半人马族小天体的发现，使小行星与彗星之间的界限变得模糊。21 世纪初，一些具备长周期彗星轨道却呈现内太阳系小行星特征的小天体被称为曼岛彗星（Manx comets）。它们仍被归类为彗星，例如 C/2014 S3（PANSTARRS）\(^\text{[6]}\)。在 2013 年至 2017 年之间，共发现了 27 颗曼岛彗星\(^\text{[7]}\)。

截至 2021 年 11 月，共已知 4,584 颗彗星\(^\text{[8]}\)。然而，这仅占可能存在的彗星总量的一小部分，因为太阳系外部（奥尔特云）中类彗星天体的储库估计约为一万亿个\(^\text{[9][10]}\)。平均每年约有一颗彗星可被肉眼看见，尽管其中许多较为暗淡、外观并不特别醒目\(^\text{[11]}\)。特别明亮的彗星被称为“伟大彗星”（great comets）。人类已利用无人探测器访问过彗星，例如美国 NASA 的“深度撞击号”（Deep Impact），其通过撞击坦普尔 1 号彗星形成陨坑以研究其内部结构；又如欧洲航天局的“罗塞塔号”（Rosetta），它成为首个在彗星上着陆的机器人探测器\(^\text{[12]}\)。
\subsection{词源}
\begin{figure}[ht]
\centering
\includegraphics[width=8cm]{./figures/6a3173d4cf8911ea.png}
\caption{公元 729 年的《盎格鲁－撒克逊编年史》以及尊敬的比德（the Venerable Bede）的著作中均提到了彗星。} \label{fig_Comet_2}
\end{figure}
单词“comet（彗星）”源自古英语 cometa，其来源为拉丁语 comēta 或 comētēs，而这些拉丁词又是对希腊语 κομήτης 的罗马化形式，意为“留着长发的（事物）”。《牛津英语词典》指出，术语（ἀστὴρ） κομήτης 在希腊语中早已表示“长发之星，即彗星”。κομήτης 源自动词 κομᾶν（koman），意为“留长发”，而该动词又派生自名词 κόμη（komē），意为“头发”，并被用来表示“彗尾”\(^\text{[15][16]}\)。

彗星的天文学符号（在 Unicode 中表示）为 U+2604 ☄ COMET，由一个小圆盘及其后三条类似发丝的延伸部分组成\(^\text{[17]}\)。
\subsection{物理特征}
\begin{figure}[ht]
\centering
\includegraphics[width=8cm]{./figures/37fee7dd52548b31.png}
\caption{彗星的结构} \label{fig_Comet_3}
\end{figure}
\subsubsection{彗核}
\begin{figure}[ht]
\centering
\includegraphics[width=8cm]{./figures/43510b130873771c.png}
\caption{在航天器飞越期间拍摄的 103P/Hartley 彗核图像。该彗核长度约为 2 千米。} \label{fig_Comet_4}
\end{figure}
彗星的固体核心结构被称为彗核。彗核由岩石、尘埃、水冰以及冻结的二氧化碳、一氧化碳、甲烷和氨等物质的混合体构成\(^\text{[18]}\)。因此，根据弗雷德·惠普尔（Fred Whipple）的模型，它们常被形象地称为“肮脏的雪球”（dirty snowballs）\(^\text{[19]}\)。尘埃含量更高的彗星则被称为“冰质尘球”（icy dirtballs）\(^\text{[20]}\)。“冰质尘球”这一术语源于 2005 年 7 月 NASA“深度撞击号”（Deep Impact）任务向坦普尔 1 号彗星（Comet 9P/Tempel 1）发送“撞击器”探测器并观测其碰撞后的结果。2014 年的一项研究表明，彗星类似于“炸冰淇淋”，即其表面由密集的晶体冰与有机化合物混合构成，而内部冰层则更冷且密度更低\(^\text{[21]}\)。

彗核表面通常干燥、多尘或岩质，表明冰层隐藏在厚度数米的表层壳体之下。彗核中含有多种有机化合物，其中可能包括甲醇、氰化氢、甲醛、乙醇、乙烷，以及可能更复杂的分子，如长链烃类与氨基酸\(^\text{[22][23]}\)。2009 年确认，NASA “星尘号”（Stardust）任务回收的彗星尘埃中发现了氨基酸甘氨酸\(^\text{[24]}\)。2011 年 8 月，一份基于 NASA 对地球陨石研究的报告提出，DNA 与 RNA 的组成单元（腺嘌呤、鸟嘌呤及相关有机分子）可能形成于小行星和彗星\(^\text{[25][26]}\)。

彗核外层表面具有极低反照率，使其成为太阳系中反射率最低的天体之一。贾托号（Giotto）探测器发现，哈雷彗星（1P/Halley）的彗核仅反射约四个百分点的入射光\(^\text{[27]}\)；而“深空一号”（Deep Space 1）发现波罗雷彗星（Comet Borrelly）的表面反射率甚至不足 3.0\%\(^\text{[27]}\)。相比之下，沥青可反射约 7\%。彗核表面的暗色物质可能由复杂的有机化合物构成。太阳加热会驱散较轻、易挥发的化合物，而留下更大的有机分子，这些物质通常极为黝黑，如焦油或原油。彗核表面的低反照率使其能够吸收驱动释气过程的热量\(^\text{[28]}\)。

已观测到的彗核半径可达 30 千米（19 英里）\(^\text{[29]}\)，但要准确测定彗核的大小并不容易\(^\text{[30]}\)。322P/SOHO 的彗核直径可能只有 100–200 米（330–660 英尺）\(^\text{[31]}\)。尽管观测仪器的灵敏度不断提高，但仍缺乏对更小彗星的探测，这使部分研究者推测直径小于 100 米（330 英尺）的彗星数量可能确实稀少\(^\text{[32]}\)。已知彗星的平均密度估计约为 $0.6 g/cm^3$（0.35 oz/cu in）\(^\text{[33]}\)。由于质量较低，彗核不会在自引力作用下成为球形，因此呈现不规则形状\(^\text{[34]}\)。
\begin{figure}[ht]
\centering
\includegraphics[width=8cm]{./figures/7e9645911c29b572.png}
\caption{彗星 81P/Wild 在光照面和阴暗面都呈现出喷流，其表面起伏鲜明，且十分干燥。} \label{fig_Comet_5}
\end{figure}
大约六个百分点的近地小行星被认为是不再发生释气作用的彗星灭绝彗核\(^\text{[35]}\)，其中包括 14827 Hypnos 与 3552 Don Quixote。

“罗塞塔号”（Rosetta）与“菲莱”（Philae）着陆器的探测结果显示，67P/丘留莫夫－格拉西缅科彗星（67P/Churyumov–Gerasimenko）的彗核不存在磁场\(^\text{[36][37]}\)，这表明磁性可能并未在早期微行星形成过程中发挥作用。此外，罗塞塔号上的 ALICE 光谱仪确定，在距离彗核约 1 千米范围内，由太阳辐射对水分子进行光电离所产生的电子，而非此前认为的来自太阳的光子，是导致彗核释放出的水与二氧化碳分子在彗发中发生分解的主要因素\(^\text{[38][39]}\)。菲莱着陆器上的仪器在彗星表面检测到至少十六种有机化合物，其中有四种（乙酰胺、丙酮、异氰酸甲酯和丙醛）是首次在彗星上被发现的\(^\text{[40][41][42]}\)。
\begin{figure}[ht]
\centering
\includegraphics[width=10cm]{./figures/8c2117a476d375ac.png}
\caption{} \label{fig_Comet_6}
\end{figure}
\subsubsection{彗发}
\begin{figure}[ht]
\centering
\includegraphics[width=8cm]{./figures/3f4fb64d5886d128.png}
\caption{Borrelly 彗星呈现喷流，但其表面不存在冰。} \label{fig_Comet_7}
\end{figure}
\begin{figure}[ht]
\centering
\includegraphics[width=8cm]{./figures/2d693476f2b0c944.png}
\caption{近日点前不久由哈勃望远镜拍摄的 ISON 彗星图像[50]} \label{fig_Comet_8}
\end{figure}
由此释放出的尘埃和气体会在彗星周围形成一个巨大而极其稀薄的大气层，称为“彗发”（coma）。太阳的辐射压与太阳风施加在彗发上的力会形成一条巨大的彗尾，且彗尾始终指向远离太阳的方向\(^\text{[51]}\)。

彗发通常由水和尘埃组成，其中水占彗星在距太阳 3 至 4 个天文单位（4.5 亿至 6 亿千米；2.8 亿至 3.7 亿英里）范围内从彗核中流出的挥发物的 90\% 之多\(^\text{[52]}\)。母体分子 $H_2O$主要通过光解作用被破坏，而光电离的贡献要小得多；与光化学过程相比，太阳风对水的破坏作用仅占很小比例\(^\text{[52]}\)。较大的尘埃颗粒会沿着彗星轨道路径残留，而较小的颗粒则被阳光压力推向远离太阳的方向进入彗尾\(^\text{[53]}\)。

虽然彗星的固体彗核通常小于 60 千米（37 英里）宽，但彗发的大小可能达到数千甚至数百万千米，有时甚至比太阳还大\(^\text{[54]}\)。例如，2007 年 10 月的一次爆发过后大约一个月，彗星 17P/Holmes 的稀薄尘埃大气层短暂地变得比太阳还大\(^\text{[55]}\)。1811 年的“大彗星”其彗发的直径大约与太阳相当\(^\text{[56]}\)。尽管彗发可以变得非常庞大，但其尺寸会在彗星经过火星轨道附近（约 1.5 AU，即 2.2 亿千米；1.4 亿英里）时缩小\(^\text{[56]}\)。在这一距离，太阳风强度足以将彗发中的气体和尘埃吹散，从而反而使彗尾变得更长\(^\text{[56]}\)。离子彗尾的长度可达到一个天文单位（1.5 亿千米）或更长\(^\text{[55]}\)。
\begin{figure}[ht]
\centering
\includegraphics[width=8cm]{./figures/f1b7d1bdbd91578d.png}
\caption{C/2006 W3（Christensen）释放碳气体（红外图像）} \label{fig_Comet_9}
\end{figure}
彗发与彗尾都由太阳照亮，并可能在彗星穿过内太阳系时变得可见；其中尘埃直接反射阳光，而气体则因电离而产生辉光\(^\text{[57]}\)。大多数彗星过于暗淡，以至于没有望远镜无法看见，但每隔数十年会有少数彗星亮到可被肉眼观测\(^\text{[58]}\)。有时彗星会经历一次巨大而突然的气体与尘埃爆发，在此期间彗发的尺寸会显著增大，持续一段时间。2007 年的霍姆斯彗星（Comet Holmes）便发生了这一现象\(^\text{[59]}\)。

1996 年，人们发现彗星能够发射 X 射线\(^\text{[60]}\)。这一发现令天文学家非常惊讶，因为 X 射线通常与极高温的天体相关。Thomas E. Cravens 在 1997 年初首次提出了解释\(^\text{[61]}\)。X 射线的产生源于彗星与太阳风之间的相互作用：当高电荷态的太阳风离子穿过彗发大气层时，它们会与彗星的原子和分子发生碰撞，并在这一过程中“夺取”一个或多个电子，这一机制称为“电荷交换”（charge exchange）。电子被太阳风离子俘获之后，会通过向基态去激发的过程发射出 X 射线和远紫外光子\(^\text{[62]}\)。
\subsubsection{弓形激波}
弓形激波是由太阳风与彗星电离层相互作用形成的，而彗星的电离层则由彗发中气体的电离所产生。随着彗星逐渐靠近太阳，释气速率不断增加，使得彗发膨胀；阳光会进一步电离彗发中的气体。当太阳风经过这一离子化的彗发区域时，弓形激波便会出现。

第一次观测是在 1980 年代与 1990 年代，当时多艘航天器先后飞掠了 21P/Giacobini–Zinner\(^\text{[63]}\)、1P/Halley\(^\text{[64]}\)与 26P/Grigg–Skjellerup\(^\text{[65]}\)彗星。结果发现，彗星的弓形激波比例如地球那样的行星弓形激波更宽且更加平缓。这些观测均发生在近日点附近，当时弓形激波已完全形成。

罗塞塔号（Rosetta）在 67P/Churyumov–Gerasimenko 彗星上观测到弓形激波的早期形成阶段，当时彗星正向太阳靠近，释气活动逐渐增强。这个尚在成长中的弓形激波被称为“婴儿弓形激波”（infant bow shock）。婴儿弓形激波是不对称的，并且相对于彗核距离而言，比完全发育的弓形激波更为宽广\(^\text{[66]}\)。
\subsubsection{彗尾}
\begin{figure}[ht]
\centering
\includegraphics[width=8cm]{./figures/ec7df3f317631c33.png}
\caption{彗星在接近太阳轨道时彗尾的典型指向方向} \label{fig_Comet_10}
\end{figure}
在外太阳系中，彗星保持冻结且处于不活跃状态，由于体积很小，从地球上极难或几乎不可能被探测到。基于哈勃空间望远镜的观测，有研究报告在柯伊伯带中统计性地探测到不活跃的彗星彗核[67][68]，但这些探测结果曾受到质疑[69][70]。当彗星进入内太阳系时，太阳辐射会使其中的挥发性物质升华并从彗核中逸出，同时携带尘埃。

逸出的尘埃与气体分别形成各自独立的彗尾，方向略有不同。尘埃彗尾通常沿彗星轨道附近被遗留，从而形成一条常常呈弯曲状的彗尾，称为 II 型或尘埃彗尾（dust tail）[57]。同时，离子彗尾（ion tail），亦称 I 型彗尾，由气体组成，始终直接指向远离太阳的方向，因为气体比尘埃更强烈地受到太阳风作用，并沿磁力线运动，而非遵循轨道轨迹[71]。在某些情况下——例如当地球经过彗星的轨道平面时——可以看到一条反尾（antitail），其方向与离子彗尾和尘埃彗尾相反[72]。
\begin{figure}[ht]
\centering
\includegraphics[width=8cm]{./figures/1dacb212df9b211e.png}
\caption{彗星示意图：展示由太阳风形成的尘埃轨迹、尘埃彗尾与离子气体彗尾} \label{fig_Comet_11}
\end{figure}
对反尾（antitail）的观测对太阳风的发现起到了重要作用[73]。离子彗尾是由于彗发中粒子被太阳紫外辐射电离而形成的。一旦粒子被电离，它们便获得净正电荷，从而在彗星周围产生一个“感应磁层”（induced magnetosphere）。彗星及其感应磁场对向外流动的太阳风粒子构成阻碍。由于彗星与太阳风之间的相对轨道速度为超音速，彗星在太阳风流向的上游会形成一个弓形激波。在这一弓形激波中，彗星离子（称为“拾取离子” pick-up ions）大量聚集，并使太阳磁场被等离子体“负载”，使磁力线在彗星周围“披覆”（drape），从而形成离子彗尾[74]。

若离子彗尾的负载程度足够大，磁力线会被压缩到某一距离处发生磁重联（magnetic reconnection）。这会导致一次“彗尾断裂事件”（tail disconnection event）[74]。这一现象已多次被观测到，其中一个著名事件发生在 2007 年 4 月 20 日：当恩克彗星（Encke’s Comet）穿过一次日冕物质抛射（CME）时，其离子彗尾完全被切断。该事件由 STEREO 空间探测器记录[75]。

2013 年，欧洲航天局（ESA）的科学家报告称，金星的电离层会向外流出，其行为与类似条件下彗星的离子彗尾相似[76][77]。
\subsubsection{喷发}
\begin{figure}[ht]
\centering
\includegraphics[width=8cm]{./figures/e97b94991a13b0fc.png}
\caption{103P/Hartley 的气体与冰雪喷流} \label{fig_Comet_12}
\end{figure}
不均匀的加热会使新生成的气体从彗核表面的薄弱区域喷出，类似于间歇泉[78]。这些气体和尘埃的喷流能够使彗核旋转，甚至将其分裂[78]。2010 年的研究显示，干冰（冻结的二氧化碳）的升华能够为从彗核喷出的物质喷流提供动力[79]。对 Hartley 2 彗星的红外成像显示了此类喷流从其表面喷出，并携带尘埃颗粒进入彗发[80]。
\subsection{轨道特征}
大多数彗星是太阳系内的小天体，其轨道为细长的椭圆轨道，轨道的一部分将它们带到靠近太阳的区域，而其余时间则位于太阳系的遥远外侧[81]。彗星通常根据其轨道周期长短进行分类：周期越长，其轨道椭圆的离心率越大。
\subsubsection{短周期}
周期彗星或短周期彗星通常被定义为那些轨道周期小于 200 年的彗星[82]。它们通常在黄道面内、并以与行星相同的方向运行[83]。其轨道在远日点时往往延伸至外行星（木星及更远）的区域；例如，哈雷彗星的远日点略微超出海王星轨道。远日点接近某一主要行星轨道的彗星被称为该行星的“族”（family）[84]。一般认为，这类彗星族是由行星将原本属于长周期彗星的天体捕获至较短轨道而形成的[85]。

在短轨道周期的极端情况下，恩克彗星（Encke's Comet）的轨道甚至未到达木星轨道，被称为恩克型彗星（Encke-type comet）。轨道周期小于 20 年、且轨道倾角对黄道面不超过 30 度的短周期彗星被称为传统木星族彗星（Jupiter-family comets, JFCs）[86][87]。类似哈雷彗星、轨道周期介于 20 至 200 年之间、且轨道倾角范围从 0 度一直到超过 90 度的，则称为哈雷型彗星（Halley-type comets, HTCs）[88][89]。截至 2025 年 7 月，已知恩克型彗星共有 74 颗（其中 6 颗被归类为近地天体 NEOs），HTCs 有 109 颗（其中 36 颗为 NEOs），而 JFCs 有 815 颗（其中 153 颗为 NEOs）[90]。

近期发现的主带彗星（main-belt comets）构成一个独立的类别，其轨道更为接近圆形，并位于小行星带内部运行[91][92]。

由于其椭圆轨道经常使它们接近巨行星，彗星会受到进一步的引力扰动[93]。短周期彗星的远日点往往与某一巨行星的半长轴相一致，其中木星族彗星（JFCs）是最大的一类[87]。来自奥尔特云的彗星在经历与巨行星的近距离遭遇后，其轨道常会受到强烈的引力影响。木星因其质量超过其他所有行星的总和，是引力扰动最显著的来源。这些扰动能够将长周期彗星偏转至较短周期的轨道上[94][95]。

基于轨道特征，短周期彗星通常被认为起源于半人马小行星（centaurs）以及柯伊伯带／散射盘（Kuiper belt/scattered disc）[96] —— 一个位于海王星外侧区域的盘状小天体群；而长周期彗星的来源则被认为是更遥远的球形奥尔特云（Oort cloud），其名称来自假设其存在的荷兰天文学家 Jan Hendrik Oort[97]。在这些遥远区域中，巨量彗星状天体被认为以近圆轨道绕太阳运行。有时外行星（对柯伊伯带天体而言）或邻近恒星（对奥尔特云天体而言）的引力会将其中某个天体抛入一条椭圆轨道，使其向太阳方向运动，从而形成可见彗星。与周期彗星的回归不同，这类新彗星的出现是不可预测的，因为它们的轨道此前从未被建立[98]。当地被抛入太阳的轨道并不断向太阳落入时，其质量会持续被剥离，而这种剥离对彗星的寿命影响巨大：剥离越严重，寿命越短，反之亦然[99]。
\subsubsection{长周期}
\begin{figure}[ht]
\centering
\includegraphics[width=8cm]{./figures/aff8f46e648a6626.png}
\caption{科霍特克彗星（红色）与地球（蓝色）的轨道示意图，展示其轨道的高离心率以及彗星在接近太阳时的高速运动。} \label{fig_Comet_13}
\end{figure}
长周期彗星具有高度偏心的轨道，其周期范围从 200 年延伸至数千年、甚至数百万年[100]。在近日点附近出现偏心率大于 1 的情况并不必然意味着彗星将离开太阳系[101]。例如，麦克诺特彗星（Comet McNaught）在 2007 年 1 月经过近日点时，其日心瞬时轨道偏心率为 1.000019，但其真实轨道仍被太阳束缚，轨道周期约为 92,600 年，因为偏心率在其远离太阳后会下降至 1 以下。要得到长周期彗星的未来轨道，必须在其离开行星区域后计算其瞬时轨道，并且应相对于太阳系质心来计算。根据定义，长周期彗星保持被太阳引力束缚；那些因与主要行星近距离掠过而被抛出太阳系的彗星，不再被认为具有“周期”。长周期彗星的远日点远远超出外行星，其轨道平面也不必与黄道面接近。诸如 C/1999 F1 与 C/2017 T2（PANSTARRS）等长周期彗星，其远日点距离可以接近 70,000 AU（0.34 pc；1.1 光年），估计轨道周期约为 600 万年。

单次显现或非周期彗星在性质上与长周期彗星类似，因为它们在内太阳系近日点附近具有近似抛物线或略呈双曲线的轨迹[100]。然而，巨行星的引力扰动会改变它们的轨道。单次显现彗星的瞬时轨道为双曲线或抛物线，使其在仅经过一次太阳后便永久离开太阳系[102]。太阳的希尔球（Hill sphere）具有约 230,000 AU（1.1 pc；3.6 光年）的不稳定最大边界[103]。只有数百颗彗星在近日点附近被观测到具有双曲线轨道（e > 1）[104]；基于日心的、未受扰动的双体拟合轨道表明，它们可能会逃离太阳系。

截至 2025 年，已有三颗偏心率显著大于 1 的天体被发现：1I/ʻOumuamua、2I/Borisov 与 3I/ATLAS，它们均被认为起源于太阳系之外。ʻOumuamua 的偏心率约为 1.2，在其于 2017 年 10 月穿过内太阳系期间并未显示出彗星活动的光学迹象，但其轨迹的变化——暗示存在释气作用——表明其很可能是一颗彗星[105]。相比之下，2I/Borisov 的偏心率约为 3.36，已被观测到具有典型的彗发结构，被认为是首次被探测到的星际彗星[106][107]。3I/ATLAS 的偏心率约为 6.1，同样具有彗发，说明它也是一颗彗星。C/1980 E1 在 1982 年近日点之前的轨道周期约为 710 万年，但其在 1980 年与木星的一次遭遇使其加速，赋予它目前已知太阳系内彗星中最大的偏心率（1.057），且观测弧度相对可靠[108]。预计不会再返回内太阳系的彗星包括 C/1980 E1、C/2000 U5、C/2001 Q4（NEAT）、C/2009 R1、C/1956 R1 以及 C/2007 F1（LONEOS）。

部分权威资料将“周期彗星”（periodic comet）用于指代所有轨道周期性的彗星（即所有短周期与长周期彗星）[109]，而另一些资料则仅将其限定为短周期彗星[100]。类似地，尽管“非周期彗星”（non-periodic comet）从字面上等同于“单次显现彗星”，但在另一种用法中，它被用于指代所有非上述狭义“周期彗星”的彗星（即轨道周期超过 200 年的所有彗星）。

早期观测确实发现了一些真正的双曲线（即非周期）轨迹，但数量不超过木星引力扰动能够解释的范围。来自星际空间的彗星其速度量级与太阳附近恒星的相对速度相似（每秒数十公里）。当这类天体进入太阳系时，它们具有正的轨道比能，从而具有正的无穷远速度（$v_{\infty}$），并呈现明显的双曲线轨迹。粗略估计，在木星轨道范围内，每世纪可能出现约四颗星际双曲线彗星，误差可达一个至两个数量级[110]。
\begin{figure}[ht]
\centering
\includegraphics[width=14.25cm]{./figures/b8f95819756446e8.png}
\caption{} \label{fig_Comet_14}
\end{figure}
\subsubsection{奥尔特云与希尔斯云}
\begin{figure}[ht]
\centering
\includegraphics[width=14.25cm]{./figures/8a76749915cbfb85.png}
\caption{被认为包围着太阳系的奥尔特云示意图，并展示了柯伊伯带和小行星带以作比较。} \label{fig_Comet_15}
\end{figure}
奥尔特云被认为占据着一片广阔的空间，其内边界大约从距离太阳 2,000 至 5,000 AU（0.03–0.08 光年）[112] 开始，一直到约 50,000 AU（0.79 光年）[88]。这一云团包围着从太阳系中部——太阳本身——一直到柯伊伯带外缘的所有天体。奥尔特云由适合形成天体的可行物质组成。太阳系行星之所以存在，是因为这些由太阳引力聚集并凝结的微行星（形成行星时遗留下来的固体碎块）。这些被困微行星形成的偏心性正是奥尔特云存在的原因[113]。一些估计将其外缘置于 100,000 至 200,000 AU（1.58–3.16 光年）之间[112]。该区域可进一步分为距离太阳 20,000–50,000 AU（0.32–0.79 光年）的球状外奥尔特云，以及距离太阳 2,000–20,000 AU（0.03–0.32 光年）的甜甜圈状内云，即希尔斯云（Hills cloud）[114]。外奥尔特云与太阳的引力束缚很弱，是落入海王星轨道以内的长周期彗星（以及可能的哈雷型彗星）的来源[88]。内奥尔特云亦被称为希尔斯云，以 1981 年提出其存在的 Jack G. Hills 命名[115]。模型预测内云中包含的彗核数量应为外云的数十倍至数百倍[115][116][117]；其被视为新彗星的潜在来源，为随着时间流逝而逐渐变得稀薄的外奥尔特云提供补给。希尔斯云解释了奥尔特云在数十亿年后仍然存在的原因[118]。
\subsubsection{系外彗星}
太阳系之外的系外彗星已被探测到，并可能在银河系中普遍存在[119]。首个被探测到的系外彗星系统是 1987 年发现的织女增二（Beta Pictoris）系统，它是一颗极年轻的 A 型主序星[120][121]。截至 2013 年，共识别出 11 个此类系外彗星系统，其探测方式依赖于彗星在靠近其母恒星时释放的大量气体所造成的吸收光谱特征[119][120]。开普勒空间望远镜在十年间负责搜索太阳系外的行星与其他天体。首例凌日系外彗星于 2018 年 2 月被一个由专业天文学家与公民科学家组成的小组发现，他们在开普勒的光变曲线中识别出了这些信号[122][123]。在开普勒望远镜于 2018 年 10 月退役后，其任务由 TESS 望远镜接替。自 TESS 发射以来，天文学家利用其光变曲线在织女增二周围观测到彗星的凌日现象[124][125]。

随着 TESS 的投入使用，天文学家能够更好地利用光谱方法区分系外彗星。新行星通常通过白光光变法探测，当行星遮挡母恒星时，会在图表读数中出现对称的下降。然而，对这些光变曲线的进一步分析显示，这些不对称下降的模式实际上是由彗尾或数百条彗尾引起的[126]。
\subsection{彗星的影响}
\begin{figure}[ht]
\centering
\includegraphics[width=8cm]{./figures/b8829ab21e3c6012.png}
\caption{英仙座流星示意图} \label{fig_Comet_16}
\end{figure}
\subsubsection{彗星与流星雨的关联}
当彗星在接近太阳的过程中受热时，其冰质组分的释气作用会释放出固体碎屑，这些碎屑太大，辐射压与太阳风无法将其完全吹散[127]。如果地球的轨道恰好穿过这条主要由细小岩质颗粒组成的碎屑轨迹，那么当地球经过其中时就可能出现流星雨。碎屑轨迹越密集，产生的流星雨越短暂但越强烈；而较稀疏的碎屑轨迹则会形成时间较长但强度较低的流星雨。通常而言，碎屑轨迹的密度与母彗星释放这些物质的时间有直接关系[128][129]。例如，英仙座流星雨每年 8 月 9 日至 13 日出现，此时地球通过斯威夫特－塔特尔彗星（Comet Swift–Tuttle）的轨道。哈雷彗星则是 10 月出现的猎户座流星雨的来源[130][131]。
\subsubsection{彗星与生命的影响}
在早期地球形成阶段，许多彗星和小行星与地球发生过碰撞。许多科学家认为，大约 40 亿年前彗星轰击年轻地球，为地球海洋带来了如今巨量的水，或者至少带来了其中相当大的一部分。然而，也有研究对这一观点提出质疑[132]。在彗星中探测到的有机分子（包括多环芳烃）[21] 的显著含量，引发了关于彗星或陨石可能将生命前体——甚至生命本身——带到地球的推测[133]。2013 年有研究提出，像彗星这样岩石表面与冰表面发生碰撞的事件具有通过冲击合成生成构成蛋白质的氨基酸的潜力[134]。彗星进入大气层的高速运动，以及初次撞击产生的巨大能量，使得小分子能够凝结并形成更大的大分子，这些大分子为生命的出现奠定了基础[135]。2015 年，科学家在 67P 彗星的释气物质中发现了大量的分子氧，这暗示该分子可能比原先认为的更为常见，从而使其作为生命迹象的意义降低[136]。

人们推测，在长时间尺度上，彗星撞击可能向月球输送了大量水，其中一部分可能以月球冰的形式保留下来[137]。彗星和流星体撞击也被认为是玻陨石（tektites）和澳洲玻陨石（australites）存在的原因[138]。
\subsubsection{对彗星的恐惧}
在 1200 至 1650 年间，欧洲人对彗星的恐惧达到高峰，将其视为上帝的行为或灾难即将来临的征兆[139]。例如，在 1618 年“大彗星”出现后的第二年，Gotthard Arthusius 发表了一份小册子，声称该彗星预示着世界末日临近[140]。他列出了十页与彗星相关的灾难，包括“地震、洪水、河道改变、冰雹、炎热干燥的天气、歉收、瘟疫、战争、背叛与物价飞涨”等[139]。

到 1700 年，大多数学者已经认识到这些灾难无论是否出现彗星都会发生。然而，威廉·惠斯顿（William Whiston）在 1711 年利用爱德蒙·哈雷记录的彗星观测数据写道，1680 年的大彗星具有 574 年的周期性，并通过向地球倾泻洪水的方式导致了《创世纪》中的全球性大洪水。他的观点再次激发了人们对彗星的恐惧，这一次不是将其视为灾难的征兆，而是视为地球的直接威胁[139]。1910 年光谱分析在哈雷彗星彗尾中发现了有毒气体氰（cyanogen）[141]，导致公众恐慌性购买防毒面具以及各种伪科学的“抗彗星药片”和“抗彗星雨伞”[142]。
\subsection{彗星的命运}
\subsubsection{离开（被抛出）太阳系}
如果一颗彗星的速度足够快，它可能会离开太阳系。这样的彗星遵循双曲线的开放轨道，因此被称为双曲线彗星。已知的太阳系内双曲线彗星只有在与太阳系中其他天体（如木星）相互作用后才会被抛出[143]。一个例子是 C/1980 E1 彗星，它在 1980 年与木星的一次近距离掠过后，其原本约 710 万年的绕日轨道被改变为双曲线轨道[144]。而诸如 1I/ʻOumuamua、2I/Borisov 和 3I/ATLAS 等星际彗星从未绕太阳运行，因此不需要第三天体的相互作用即可离开太阳系。
\subsubsection{灭绝}
木星族彗星（JFCs）与长周期彗星表现出非常不同的亮度衰减规律。木星族彗星的活跃寿命约为 10,000 年，或大约 1,000 个轨道周期，而长周期彗星的衰减速度要快得多。只有 10\% 的长周期彗星能在多于 50 次的小近日点掠过中存活，而只有 1\% 能存活超过 2,000 次掠过[35]。最终，彗核中的大部分挥发性物质都会蒸发，彗星会变成一个小型、黑暗、惰性的岩石或碎块，其外观可能类似小行星[145]。一些椭圆轨道的小行星如今被识别为灭绝彗星[146][147][148][149]。大约六个百分点的近地小行星被认为是灭绝的彗星彗核[35]。
\subsubsection{破裂与碰撞}
某些彗星的彗核可能十分脆弱，这一结论得到了彗星分裂现象的支持[150]。一次重大彗星破裂事件是 1993 年发现的舒梅克－列维 9 号彗星（Comet Shoemaker–Levy 9）。其在 1992 年 7 月与木星的一次近距离遭遇中被撕裂成多个碎片，并在 1994 年 7 月的六天时间里，这些碎片陆续坠入木星大气层——这是天文学家首次观测到太阳系内两个天体之间的碰撞[151][152]。其他发生分裂的彗星包括 1846 年的 3D/Biela 以及 1995 至 2006 年间的 73P/Schwassmann–Wachmann[153]。希腊历史学家埃福鲁斯（Ephorus）报告称，早在公元前 372–373 年的冬季就曾有彗星分裂[154]。人们认为彗星可能因热应力、内部气体压力或撞击而发生分裂[155]。

42P/Neujmin 和 53P/Van Biesbroeck 彗星似乎是同一母体彗星的碎片。数值积分显示，这两颗彗星在 1850 年 1 月曾非常接近木星，且在 1850 年之前，它们的轨道几乎完全相同[156]。另一个由分裂事件产生的彗星族是 Liller 彗星家族，包括 C/1988 A1（Liller）、C/1996 Q1（Tabur）、C/2015 F3（SWAN）、C/2019 Y1（ATLAS）以及 C/2023 V5（Leonard）[157][158]。

一些彗星在近日点通过期间发生了破裂，包括著名的大彗星 West 和 Ikeya–Seki。比埃拉彗星（Biela's Comet）是一个典型例子：它在 1846 年穿过近日点时分裂为两个部分。这两个部分在 1852 年仍被分开观测到，但此后再也未被看到。相反，在 1872 年和 1885 年（彗星本应出现的时段）出现了壮观的流星雨。每年 11 月出现的小型流星雨——仙女座流星雨（Andromedids）——便是在地球穿过比埃拉彗星轨道时产生的[159]。

有些彗星则以更加戏剧性的方式终结——要么坠入太阳[160]，要么与行星或其他天体发生碰撞。在早期太阳系中，彗星与行星或卫星的碰撞十分常见：例如月球上的许多陨石坑可能是彗星造成的。最近一次彗星与行星的碰撞发生在 1994 年 7 月，当时舒梅克－列维 9 号彗星分裂成碎片并与木星相撞[161]。
\begin{figure}[ht]
\centering
\includegraphics[width=14.25cm]{./figures/fad311c089faff94.png}
\caption{} \label{fig_Comet_17}
\end{figure}

